
\documentclass[a4paper,twoside,openany]{article}

% --- Langue / encodage ---
\usepackage[utf8]{inputenc}
\usepackage[T1]{fontenc}
\usepackage[french]{babel}
\usepackage{lmodern}
\usepackage{microtype}
\sloppy


\title{1355 et 1360 Concession d’une rente par le comte de Savoie Amedé VI à Jean de Beaufort sur l’alpéage et la leyde de la châtellenie de Beaufort
	\\(Archives ????)}



\begin{document}
\maketitle

\section{Présentation du document}

Le document suivant, daté du 25 juin 1360,  éclaire l'ancienneté des droits seigneuriaux et pastoraux attachés aux montagnes de la châtellenie de Beaufort. Le comte Amédée VI de Savoie y assigne à Jean de Beaufort ne rente de dix florins d'or sur la leyde et l'alpeage des montagnes de Beaufort, tout en réservant le fief, le domaine direct et l'hommage.

\subsection{Texte latin}

Nos Amedeus, comes Sabaudie, notum facimus universis
quod, cum, consideratione servitiorum nobis fideliter
impensorum per dictos fideles nostros Dominum Johanen
(de Belloforti) militem, Guigonem, Petrum domicillos et
Petrum Donzelli de eodem loco, in habenda et obtinenda
possessione terre nostre Faucigniaci, pro nobis et nostris
successoribus dederimus et concesserimus eisdem in
augmentum feudorum que tenent a nobis et sub eisdem
homagiis ad que pro suis aliis feudis tenebantur cuilibet
ipsorum videlicet decem florenos auri de redditu per annum
pro se et successoribus cujuslibet ipsorum in Castellania
nostra Bellifortis assignandos prout in nostris litteris dicte
donationis, datis Gebenn. die vicesima prima Julii anno
millesimo trecentesimo quinquagesimo quinto plenius
videtur contineri, de quibus nondum est assignatio facta.
Volentes observare premissa, dicti Domini Johanis militis
supplicationi inclinati, super hiis habita relatione Domini
Aymonis de Chaland et Guillelmi Boni quibus hoc
commisimus, referentium assignationem infrasciptam
utilem esse.
pro nobis et nostris successoribus, dicto Domino Johanni pro
se et suis successoribus et heredibus universis, in
assignationem decem florenorum auri dictorum, damus,
tradimus et assignamus leyda[m] et alpagium nostrum
montium castellanie Bellifortis pro quibus levatur annis
singulis in festo Sancti Christophori, a quolibet faciente
fructus in fructus montium usui diei cujus medietas ad nos in
solidum pro leyda et due partes alterius medietatis pro
alpagio et tertia pars certis nobilibus pertinent noscuntur de
quibus pro nostris juribus predictis septem quintalia cum
dimidio caseorum nobis computantur ad presens, quam
quidem leydam et alpagium cum ipsorum emolumentis
existentibus et juribus quibus cumque dicto Domino Johanni
pro se et suis ut supra pro assignatione decem florenorum,
in augmentum feudorum que tenet a nobis et sub eodem
homagio… perpetuo tradimus atque damus et ipsum
recipientem pro se et suis… per presentes investimus sibi
concedentes omnia jura et actiones nobis competentes in
predictis et singulis ipsorum… nihil retinentes preter
feudum, directum dominium et homagium predict. Nos ea
possidere constituentes ipsius Johannis procuratorem
nomine donec ipsorum corporalem possessionem accipere
et retinere valeat… (suit la formule de garantie ordinaire).
Datum Camberiaci, die 25a Junii, anno 1360… sub nostro
sigillo in testimonium premissorum per Dominum relatione
Dominorum Aymonis de Challand, Petri de Montegil ?
Johannis Ravisiis, cancellario, Guillelmi Boni facta per dictum
Amonem de Challand…

\subsection{Traduction française de travail}


Nous, Amédée, comte de Savoie, faisons savoir à tous que,
considérant les services qui nous ont été fidèlement rendus
par nos fidèles : messire Jean (de Beaufort), chevalier,
Guigon, Pierre, damoiseau, et Pierre Donzel, du même lieu,
dans l’acquisition et la conservation de la possession de
notre terre du Faucigny, nous avons donné et concédé à
chacun d’eux, pour nous et nos successeurs, en
augmentation des fiefs qu’ils tiennent de nous, et sous les
mêmes hommages auxquels ils sont tenus pour leurs autres
fiefs : dix florins d’or de revenu par an, pour chacun d’eux et
pour leurs successeurs, revenu qui devra être assigné dans
notre châtellenie de Beaufort, ainsi qu’il est plus pleinement
contenu dans nos lettres de ladite donation, données à
Genève le 21 juillet 1355. Toutefois, l’assignation de ce
revenu n’a pas encore été effectuée.
Voulant observer ce qui précède, et inclinés à la requête du
dit messire Jean, chevalier, après avoir reçu un rapport à ce
sujet de messire Aymon de Chaland et Guillaume Bon,
auxquels nous avions confié cette mission, ceux-ci
rapportant que l’assignation ci-dessous était utile.
Pour nous et nos successeurs, au dit messire Jean, pour lui et
pour tous ses successeurs et héritiers, en assignation des dix
florins d’or susdits, nous donnons, livrons et assignons la
leyde et l’alpage qui nous appartiennent dans les montagnes
de la châtellenie de Beaufort.
Ces revenus sont levés chaque année à la fête de saint
Christophe, sur tous ceux qui tirent profit des montagnes. Le
produit de ces revenus est réparti ainsi : la moitié nous
revient intégralement au titre de la leyde, les deux tiers de
l’autre moitié reviennent pour l’alpage, le dernier tiers
appartient à certains nobles. De ces droits qui nous
appartiennent, sept quintaux et demi de fromages nous sont
actuellement comptés. Cette leyde et cet alpage, avec tous
les revenus et droits qui y sont attachés, nous les donnons et
les concédons à perpétuité au dit messire Jean, pour lui et
les siens, pour tenir lieu de l’assignation des dix florins, en
augmentation des fiefs qu’il tient de nous et sous le même
hommage. Par les présentes nous l’en investissons, lui
concédant tous les droits et actions qui nous appartiennent
sur ces biens, ne réservant pour nous que le fief, la
seigneurie directe et l’hommage. Et nous constituons pour le
moment le procureur dudit Jean pour posséder ces biens en
son nom, jusqu’à ce qu’il puisse en prendre possession
corporelle effective. (suit la clause de garantie ordinaire).
Donné à Chambéry, le 25 juin 1360, sous notre sceau en
témoignage de ce qui précède, sur rapport des seigneurs
Aymon de Challand, Pierre de Montegil ?, Jean Ravisii,
chancelier, et Guillaume Bon, rapport fait par ledit Aymon de
Challand.
(Tiré des archives du Palais de l’Isle à Annecy, le 2 juin 1693,
par Nicollin, clavaire, et trouvé manuscrit, sans numéro, aux
archives de Villard-Chabod.)

\end{document}
